\section*{अध्याय \ref{ch:g10:why-clouds-form} --- बादल क्यों बनते हैं}

\begin{solution}{exo:g10:why-clouds-form:1}
\omterm{def:g10:why-clouds-form:humidity}{सापेक्ष आर्द्रता} हवा की जलवाष्प की मात्रा है, जिसे उसी तापमान और दाब पर उसकी
\omterm{def:g10:why-clouds-form:humidity}{संतृप्ति} मात्रा से भाग दिया गया हो, प्रतिशत में व्यक्त।
\omterm{def:g10:why-clouds-form:dew-point}{ओसांक} वह तापमान है जिस तक हवा को, उसकी जल-मात्रा और दाब बनाए रखते हुए,
ठंडा करना पड़ता है ताकि वह \omterm{def:g10:why-clouds-form:humidity}{संतृप्त} हो जाए।
\end{solution}

\begin{solution}{exo:g10:why-clouds-form:2}
\num{3.8}, \num{7.6}, \num{14.7} और \qty{27.2}{g/kg}। हर
\qty{10}{\degreeCelsius} इसे लगभग 2 से गुणा करता है ($7.6/3.8 = 2.0$,
$14.7/7.6 \approx 1.9$, $27.2/14.7 \approx 1.9$)।
\end{solution}

\begin{solution}{exo:g10:why-clouds-form:3}
\qty{24}{\degreeCelsius} पर \omterm{def:g10:why-clouds-form:humidity}{संतृप्ति} मात्रा: \qty{18.9}{g/kg}; \omterm{def:g10:why-clouds-form:humidity}{सापेक्ष
आर्द्रता} $10.0/18.9 \approx \qty{53}{\percent}$। \omterm{def:g10:why-clouds-form:humidity}{संतृप्ति} मात्रा
\qty{14}{\degreeCelsius} पर \qty{10.0}{g/kg} है: यही \omterm{def:g10:why-clouds-form:dew-point}{ओसांक} है।
\end{solution}

\begin{solution}{exo:g10:why-clouds-form:4}
गर्मियों में हवा नम होती है और उसका \omterm{def:g10:why-clouds-form:dew-point}{ओसांक} गिलास के तापमान से काफ़ी ऊपर होता है:
गिलास को छूने वाली हवा अपने \omterm{def:g10:why-clouds-form:dew-point}{ओसांक} से नीचे ठंडी हो जाती है और उस पर जल संघनित होता
है। सर्दियों में घर के भीतर की बहुत शुष्क हवा का \omterm{def:g10:why-clouds-form:dew-point}{ओसांक} गिलास के तापमान से कम होता है:
उससे सटकर ठंडी होने पर भी हवा
\omterm{def:g10:why-clouds-form:humidity}{संतृप्त} नहीं होती।
\end{solution}

\begin{solution}{exo:g10:why-clouds-form:5}
नीचे से गरम होना (धूप से तपी ज़मीन के ऊपर दोपहर के कपासी मेघ); पर्वत (पवन
किसी श्रेणी पर ऊपर चढ़ने को विवश); वाताग्र (ठंडी हवा के ऊपर चढ़ती गरम हवा); अभिसरण
(भूमध्य रेखा के पास या किसी अवदाब में एक साथ बहकर मिलती पवनें)।
\end{solution}

\begin{solution}{exo:g10:why-clouds-form:6}
$125 \times (30 - 18) = \qty{1500}{m}$; $125 \times (22 - 20) =
\qty{250}{m}$। दूसरा दिन: बादल ज़मीन से केवल \qty{250}{m} ऊपर
बनते हैं।
\end{solution}

\begin{solution}{exo:g10:why-clouds-form:7}
\qty{1}{km} पर: $18 - 9.8 = \qty{8.2}{\degreeCelsius}$; \qty{2}{km} पर:
$18 - 19.6 = \qty{-1.6}{\degreeCelsius}$। नीचे बैठते हुए वह उसी दर से गरम होता है
और फिर से \qty{18}{\degreeCelsius} पर ज़मीन तक पहुँचता है।
\end{solution}

\begin{solution}{exo:g10:why-clouds-form:8}
\qty{4}{\degreeCelsius/km} के साथ असंतृप्त पार्सल
(\qty{9.8}{\degreeCelsius/km}) और \omterm{def:g10:why-clouds-form:humidity}{संतृप्त} पार्सल
(\qty{5}{\degreeCelsius/km}) दोनों अपने परिवेश से तेज़ी से ठंडे होते हैं: दोनों
स्थितियों में \omterm{def:g10:why-clouds-form:stability}{स्थिर}। \qty{12}{\degreeCelsius/km} के साथ असंतृप्त बुलबुला भी
अपने परिवेश से धीरे ठंडा होता है: वह अधिक गरम रहता है और ऊपर उठता रहता है,
और गरम हवा के बुलबुले बेरोक-टोक ऊपर उठते हैं।
\end{solution}

\begin{solution}{exo:g10:why-clouds-form:9}
$14.7 - 7.6 = \qty{7.1}{g}$ प्रति किलोग्राम हवा। \omterm{def:g10:earth-energy-budget:latent-heat}{गुप्त ऊष्मा} उस हवा को गरम करती
है जिसमें जल संघनित होता है, और उसके ठंडा होने को धीमा कर देती है।
\end{solution}

\begin{solution}{exo:g10:why-clouds-form:10}
आयतन व्यास के घन के अनुसार बढ़ते हैं: फुहार की एक बूँद के लिए $(0.5/0.02)^3 = 25^3 = 15\,625$
बूँदें; वर्षा की एक बूँद के लिए $(2/0.02)^3 = 100^3 = 1\,000\,000$।
\end{solution}

\begin{solution}{exo:g10:why-clouds-form:11}
साफ़ आकाश के नीचे ज़मीन अवरक्त विकिरित करके ठंडी होती है, और उसे छूने वाली हवा भी।
यह ठंडी हवा ऊपर की हवा से अधिक घनी होती है; ढलानों पर वह नीचे बहती है और घाटी के
तल में जमा हो जाती है। वहाँ वह रात भर ठंडी होती रहती है जब तक अपने \omterm{def:g10:why-clouds-form:dew-point}{ओसांक} तक न पहुँच
जाए, और बूँदें हवा में ही बनती हैं: कोहरा। पहाड़ियों की चोटियों पर हवा उतनी ठंडी नहीं
होती और अपने
\omterm{def:g10:why-clouds-form:dew-point}{ओसांक} से ऊपर रहती है।
\end{solution}

\begin{solution}{exo:g10:why-clouds-form:12}
ऊँचाई $z$ पर तापमान $T - a z$ है और \omterm{def:g10:why-clouds-form:dew-point}{ओसांक} $T_d - b z$;
ये बराबर होते हैं जब $(a - b) z = T - T_d$, इसलिए $z = (T - T_d)/(a - b)$।
$b = 1.7$ के साथ: $a - b = 8.1$, और $1000/8.1 \approx \qty{123}{m}$ प्रति अंश।
\end{solution}

\begin{solution}{exo:g10:why-clouds-form:13}
\omterm{def:g10:why-clouds-form:condensation-level}{संघनन तल}: $125 \times (15 - 11) = \qty{500}{m}$; वहाँ तापमान:
$15 - 9.8 \times 0.5 \approx \qty{10.1}{\degreeCelsius}$। शिखर: $10.1 - 5
\times 1.5 = \qty{2.6}{\degreeCelsius}$। शुष्क दर से \qty{2}{km} नीचे:
$2.6 + 9.8 \times 2 \approx \qty{22.2}{\degreeCelsius}$, पवनाभिमुख ओर की तुलना में लगभग सात अंश
अधिक गरम।
\end{solution}

\begin{solution}{exo:g10:why-clouds-form:14}
\qty{-20}{\degreeCelsius} पर: $126/103 \approx 1.22$, बर्फ़ के लिए \qty{22}{\percent}
अतिसंतृप्त; \qty{-10}{\degreeCelsius} पर: $287/260 \approx
1.10$, \qty{10}{\percent}। इन तापमानों में बादल जितना ठंडा होगा, बर्फ़ का लाभ
उतना बड़ा होगा: प्रक्रम बादल के उस ठंडे ऊपरी भाग में सबसे अच्छा काम करता है जहाँ
बूँदें और क्रिस्टल मिले-जुले होते हैं।
\end{solution}

\begin{solution}{exo:g10:why-clouds-form:15}
असंतृप्त: पार्सल \qty{9.8}{\degreeCelsius} ठंडा होता है, उसका परिवेश
\qty{6.5}{\degreeCelsius}: वह \qty{3.3}{\degreeCelsius} अधिक ठंडा है और
वापस नीचे बैठ जाता है। \omterm{def:g10:why-clouds-form:humidity}{संतृप्त}: वह \qty{5}{\degreeCelsius} ठंडा होता है, इसलिए वह
\qty{1.5}{\degreeCelsius} अधिक गरम है और ऊपर उठता रहता है। वायु केवल एक शर्त पर
\omterm{def:g10:why-clouds-form:stability}{अस्थिर} है: कि ऊपर उठती हवा \omterm{def:g10:why-clouds-form:humidity}{संतृप्त} हो।
\end{solution}

\begin{solution}{pb:g10:why-clouds-form:1}
\textbf{1.} \qty{14.7}{g/kg}।

\textbf{2.} $11.4/14.7 \approx \qty{78}{\percent}$।

\textbf{3.} \qty{16}{\degreeCelsius}, जहाँ \omterm{def:g10:why-clouds-form:humidity}{संतृप्ति} मात्रा
\qty{11.4}{g/kg} है।

\textbf{4.} $11.4/24.1 \approx \qty{47}{\percent}$।

\textbf{5.} नहीं: उसमें वही \qty{11.4}{g/kg} है। उसकी \omterm{def:g10:why-clouds-form:humidity}{सापेक्ष आर्द्रता}
इसलिए घटी कि अधिक गरम होकर वह अधिक धारण कर सकती थी।

\textbf{6.} उसके \omterm{def:g10:why-clouds-form:dew-point}{ओसांक} पर, \qty{16}{\degreeCelsius}।

\textbf{7.} $11.4 - 8.7 = \qty{2.7}{g}$ प्रति किलोग्राम हवा।

\textbf{8.} साफ़ रात के आकाश के नीचे घास और छतें अवरक्त विकिरित करती हैं और
हवा से अधिक ठंडी हो जाती हैं; उन्हें छूने वाली हवा बाक़ी हवा से पहले अपने \omterm{def:g10:why-clouds-form:dew-point}{ओसांक}
से नीचे ठंडी हो जाती है।

\textbf{9.} $(28.0 - 13.4)/1.5 \approx \qty{9.7}{\degreeCelsius/km}$,
व्यवहार में शुष्क दर: दोपहर के बुलबुलों ने इस परत को मिला दिया है।

\textbf{10.} $(13.4 - 0.4)/2.0 = \qty{6.5}{\degreeCelsius/km}$।

\textbf{11.} असंतृप्त पार्सल के लिए \omterm{def:g10:why-clouds-form:stability}{स्थिर}, जो अधिक तेज़ी से ठंडा होता है
(\qty{9.8}{\degreeCelsius/km}); \omterm{def:g10:why-clouds-form:humidity}{संतृप्त} पार्सल के लिए \omterm{def:g10:why-clouds-form:stability}{अस्थिर}, जो
अधिक धीरे ठंडा होता है (\qty{5}{\degreeCelsius/km})।

\textbf{12.} तापमान ऊँचाई के साथ बढ़ता है, \num{-0.3} से
\qty{1.2}{\degreeCelsius} तक: एक बहुत \omterm{def:g10:why-clouds-form:stability}{स्थिर} परत, जो ढक्कन का काम करती है।
ऊपर उठता बादल वहीं रुक जाता है: उसके शीर्ष लगभग \qty{4}{km} पर हैं।

\textbf{13.} $28 - 16 = \qty{12}{\degreeCelsius}$।

\textbf{14.} तापमान $28 - 9.8 = \qty{18.2}{\degreeCelsius}$; \omterm{def:g10:why-clouds-form:dew-point}{ओसांक}
$16 - 1.8 = \qty{14.2}{\degreeCelsius}$।

\textbf{15.} 12 से $18.2 - 14.2 = 4.0$ तक: यह
\qty{8}{\degreeCelsius} घटा, यानी \qty{8}{\degreeCelsius} प्रति किलोमीटर की दर।

\textbf{16.} एक अंश का अंतर $1/8$ km में मिट जाता है, यानी
\qty{125}{m} में।

\textbf{17.} $125 \times (20 - 16) = \qty{500}{m}$।

\textbf{18.} \omterm{def:g10:why-clouds-form:dew-point}{ओसांक} \qty{16}{\degreeCelsius} पर ही रहता है जबकि ज़मीन के पास की हवा
गरम होती जाती है: अंतर बढ़ता है, और उसके साथ वह ऊँचाई भी जहाँ
यह अंतर मिटता है।

\textbf{19.} $28 - 9.8\,z = 16 - 1.8\,z$।

\textbf{20.} $8\,z = 12$, इसलिए $z = 1.5$: दोपहर के कपासी मेघ का आधार
मैदान से \textbf{लगभग \qty{1.5}{km}} ऊपर है (जल के गुणधर्मों के साथ सटीक परिकलन
\qty{1.51}{km} देता है)।
\end{solution}
